% starter-types-ar.tex — Arabic edition of ch. 6, sourdough starter types
%
% Original: The Sourdough Framework by Hendrik Kleinwächter,
% the-sourdough-framework.com, CC BY-SA 4.0. The Arabic edition inherits
% CC BY-SA 4.0.
%
% Phase 28R R2 unit 8 reconstruction. Body text recovered verbatim from
% the R0 recovery reference (recovery-r0/text/06-starter-types.txt,
% mined from the Phase 23 website build of the accepted Arabic sources)
% with the accepted edition book-ar-fixed.pdf (b78c2a28b574288801d8000a7249af3a)
% as the fidelity authority; no text retranslated from the English.
% Figure/table/flowchart labels recovered from the accepted edition's
% renders. Row 3 of جدول 6.1 reads جميع أنواع الدقيق per the accepted
% edition (the recovery reference's القمح is a recovery artifact).
%
% \ref{} calls are wrapped in \LR{}: hyperref stores the number as
% \LRE{1}.\LRE{5} fragments (the sty's \arabicnumber \@ensure@LTR
% patch), and in an RTL line the bare dot between the two LTR islands
% reverses the fragment order, so an unwrapped \ref renders "5.1"
% where the accepted edition renders the identical construct ("6.4")
% correctly. Probe-verified 2026-11-15: \LR{\LRE{1}.\LRE{5}} renders
% "1.5". Labels stay verbatim; only the direction wrapper is added.

\chapter{أنواع بادئات الخبز الحامض}\label{ch:starter-types}

\begin{quoting}
ستلقي في هذا الفصل من الكتاب نظرة أقرب على أنواع بادئات الخبز الحامض المختلفة، وعلى صفات كل نوع وطرق استخدامه. تتميّز معظمها بمستوى الترطيب فيها، وهو ما يوفّر موازنةً بين الحمضية والانتفاخ في الحجم ومستوى الغلوتين في دقيقك.
\end{quoting}

\section{المقدمة}\label{sec:starter-types-intro}

تختلف أنواع البادئة باختلاف الدقيق المتوفر لديك. فكلما زاد النشاط البكتيري زاد ما تستهلكه الكائنات الدقيقة من غلوتين. لذا إن أردت خبز رغيف قائم بذاته \textenglish{(freestanding loaf)} فستحتاج إلى دقيق أعلى في الغلوتين. وكلما زد الغلوتين زاد ما يمكن تفكيكه منه مع بقاء تماسك العجين. وإذا كنت تعيش في بلد لا يناسب مناخه زراعة القمح ولا يوفر إلا دقيقًا أضعف، فقد يُنصح باستخدام بادئة الخبز الحامض الصلبة. فالبادئة الصلبة تعزز نشاط الخميرة وتخفف النشاط البكتيري. أما إن كنت من عشاق الخبز شديد الحمضية ويتوفر لديك قمح قوي جدًا، فيمكنك تجربة بادئة الخبز الحامض السائلة \textenglish{(liquid starter)}. والفرق الجوهري بين جميع البادئات هو كمية الماء المستخدمة فيها. البادئة العادية \textenglish{(regular starter)} تحافظ على علاقة \LR{1:1} بين الدقيق والماء، والبادئة السائلة على نسبة ماء إلى دقيق قدرها \LR{5:1}، أما البادئة الصلبة فتستخدم نصف كمية الماء مقابل الدقيق، كما يلخص الجدول~\LR{\ref{tab:starter-types-comparison}}.

\begin{table}[!htb]
\centering
% table-starter-types-ar.tex — Arabic edition of upstream
% tables/table-starter-types.tex (ch. 6, جدول 6.1)
%
% Original: The Sourdough Framework by Hendrik Kleinwächter,
% the-sourdough-framework.com, CC BY-SA 4.0. The Arabic edition inherits
% CC BY-SA 4.0.
%
% RTL mirror per the settled table mechanics: the column order is
% reversed in the declared LTR tabular (upstream {@{}lclll@{}} becomes
% {@{}lllcl@{}}) so نوع البادئة renders rightmost and البكتيريا
% leftmost; the \insert@column TRT patch sets each cell's text
% direction. The \multicolumn{2}{c} Activity span and its \cmidrule
% move with the yeast/bacterial columns they cover (declared cols
% 1--2 here). Row 3 flour-type cell reads جميع أنواع الدقيق per the
% accepted edition (the R0 recovery reference's القمح is a recovery
% artifact; the accepted edition is the fidelity authority).

\begin{tabular}{@{}lllcl@{}}
\toprule
\multicolumn{2}{c}{\thead{النشاط}} & & & \\
\cmidrule(ll){1-2}
\thead{البكتيريا} & \thead{الخميرة} & \thead{نوع الدقيق} & \thead{الترطيب \LR{(\%)}} & \thead{نوع البادئة} \\
\midrule
متوازن & متوازن & قمح قوي & \LR{\num{100}} & عادية \\
مرتفع & ضئيل & قمح قوي جدًا & \LR{\num{500}} & سائلة \\
منخفض & مرتفع & جميع أنواع الدقيق & \LR{\numrange{50}{60}} & صلبة \\
\bottomrule
\end{tabular}

\caption[أنواع مختلفة من البادئة]{مقارنة بين أنواع بادئات الخبز الحامض وخواص كل منها. الفرق الوحيد بينها هو كمية الماء (الترطيب) المستخدمة عند تغذية البادئة.}\label{tab:starter-types-comparison}
\end{table}

يمكنك تغيير نوع بادئتك بمجرد تعديل نسبة التغذية، أي كميات الدقيق والماء المستخدمة. فأنا كثيرًا ما أحوّل بادئتي من عادية إلى سائلة ثم أعيدها إلى صلبة. وبعد تغيير بيئة الكائنات الدقيقة، طبّق التغذية بالنسبة نفسها على مدى بضعة أيام حتى تتأقلم مع البيئة الجديدة. ألاحظ تغيرات عادةً بعد تغذية واحدة، لكنني أنصح بتغذيتين إلى ثلاث، تغذية واحدة يوميًا، لتحصل على أثر أقوى.

\begin{figure}[!htb]
\centering
\includegraphics[width=\textwidth]{sourdough-starter-types}
\caption[بادئة سائلة وعادية وصلبة]{ثلاثة أنواع مختلفة من البادئة جنبًا إلى جنب. لاحظ كيف تكون البادئة السائلة مغمورة في الماء؛ فترطيبها \LR{\qty{500}{\percent}} أو أكثر. أما البادئة العادية فيبلغ ترطيبها نحو \LR{\qty{100}{\percent}}، والبادئة الصلبة نحو \LR{\qtyrange{50}{60}{\percent}}.}\label{fig:starter-types}
\end{figure}

عجينك ليس في جوهره سوى بادئة خبز حامض كبيرة. فبادئتك ستتكيف وتنمو من جديد داخل عجينك الرئيسي. لكن يمكنك التأثير في الصفات التي تحملها بادئتك إلى عجينك. فكلما زاد التخمير البكتيري كان تخمير عجينك البكتيري أعلى قليلًا، وكلما زاد التخمير الخميري كان تخمير عجينك الرئيسي خميريًا أعلى قليلًا. وهذا مهم عند العمل مع بادئة ناضجة غير مغذّاة.

لنفترض أن آخر تغذية لبادئتك كانت قبل \num{48} ساعة. الأرجح حينها أن بكتيرياها نشطة جدًا بينما خميرتها شبه ساكنة. في مثل هذه الحالة يمكنك الاستغناء عن تغذية البادئة قبل صنع عجين آخر؛ استخدم فقط كمية ضئيلة جدًا من البادئة. فلكل \LR{\qty{1}{\kg}} من الدقيق آخذ نحو \LR{\qty{10}{\gram}} من البادئة (\LR{\qty{1}{\percent}} بمنظومة الحساب الخبازي). وإن كانت بادئتي فتيّة جدًا وقد غذّيتها قبل \num{6} إلى \num{8} ساعات فقط، فقد أرتقي إلى \LR{\qty{20}{\percent}} من البادئة. وكما ذكرنا سابقًا، تذكّر أن عجينك ليس شيئًا سوى بادئة كبيرة؛ سيساعدك هذا كثيرًا على تحديد خطوتك التالية الأمثل.

عند استخدام معدل تلقيح منخفض كهذا (\LR{\qty{1}{\percent}}) ستحتاج إلى دقيق أقوى عند صنع العجائن القمحية. فدقيقك يتفكك طبيعيًا بفعل النشاط الأنزيمي، وقد يستغرق إعادة نمو البادئة داخل عجين الخبز نحو \num{24} ساعة. وفي الوقت نفسه قد يكون النشاط الأنزيمي قد أضعف غلوتينك تدهورًا كبيرًا. وهذا مقبول في البادئة، لكن عجينك القمحي سوف يهبط أثناء الخبز ويفقد صفاته المعتادة (بنية لبّ هشّة). لذا يُنصح بدقيق أقوى أعلى في الغلوتين؛ إذ يتيح تخميرًا أطول قبل أن يتفكك معظم الغلوتين.

\section{البادئة العادية}\label{sec:regular-starter}

\begin{figure}[!htb]
\centering
\includegraphics[width=0.5\textwidth]{sourdough-starter}
\caption[البادئة العادية]{بادئة خبز حامض عادية بترطيب \LR{\qty{100}{\percent}} مغذّاة بدقيق الجاودار.}\label{fig:regular-sourdough-starter}
\end{figure}

تُصنع البادئة العادية بترطيب يقارب \LR{\qty{100}{\percent}}، أي أن البادئة تتكون من كميتين متساويتين من الدقيق والماء. هذه البادئة هي الأكثر شيوعًا وتعليمًا على الإطلاق، وتوازن جيد بين الخميرة والبكتيريا. بعد كل تغذية يتضخم حجم العجين تضخمًا كبيرًا، ثم يبدأ بالانهيار من جديد بعد بلوغ ذروة معينة.

أفضل طريقة للحكم على جاهزية البادئة هي مراقبة علامات مثل جيوب الهواء على حواف وعائك. استخدم أنفك أيضًا لتقييم رائحة بادئتك. وإذا شعرت أن البادئة لا تعمل بالطريقة التي تُريدها، فالأرجح أن توازن الخميرة والبكتيريا فيها مضطرب. في هذه الحالة ستساعد التغذيات اليومية المتكررة بنسبة \LR{1:5:5} (بادئة: دقيق: ماء).

البادئة العادية خيار مثالي عند استخدام دقيق قمح أو سبلت أقوى، كما تعمل جيدًا مع الجاودار والإممر \textenglish{(emmer)} والقمح الأحادي \textenglish{(einkorn)}. أما إن كان لديك دقيق ضعيف قليل الغلوتين فقط فقد تسبب هذه البادئة مشكلات؛ فبسبب النشاط البكتيري الوفير فيها سيتفكك الغلوتين بسرعة. وعند استخدامها اعتمد نحو \LR{\qtyrange{1}{20}{\percent}} من البادئة بالنسبة إلى دقيق عجينك.

بحسب البكتيريا المروّضة فيها، فإن للبادئة العادية نكهةً لبنية (حمض اللبن) أو خلّية (حمض الخليك) أو مزيجًا من النكهتين. ويمكنك تعديل نكهة بادئتك بتغيير نوعها إلى بادئة سائلة.

\section{البادئة السائلة}\label{sec:liquid-starter}

\begin{figure}[!htb]
\centering
\includegraphics[width=0.5\textwidth]{sourdough-starter-liquid}
\caption[البادئة السائلة]{تتميز بادئة الخبز الحامض السائلة بكمية ماء مرتفعة. هذه الكمية الكبيرة من الماء تعزز البكتيريا المنتجة لحمض اللبن. بعد فترة يبدأ السائل والدقيق بالانفصال عن بعضهما، وتدل الفقاعات على جانب الدقيق على أن البادئة جاهزة للاستخدام.}\label{fig:liquid-sourdough-starter}
\end{figure}

\begin{flowchart}[!htb]
\centering
% fig-liquid-starter-conversion-ar.tex — Arabic edition of the liquid
% starter conversion flowchart (ch. 6, المخطط الانسيابي 6.1)
%
% Original: The Sourdough Framework by Hendrik Kleinwächter,
% the-sourdough-framework.com, CC BY-SA 4.0. The Arabic edition inherits
% CC BY-SA 4.0.
%
% R1-FIGURES route: re-derived from the upstream post-#492 source
% figures/fig-liquid-starter-conversion.tex (node styles from
% figures/flowcharts_tikz.tex and the node/edge geometry verbatim,
% including the ++(2.2,0) |- return edge, the above=2pt Yes label and
% the (90:420:1cm) repeat arc), with the label text recovered from the
% accepted edition's figure renders (label recovery, not retranslation;
% the accepted edition book-ar-fixed.pdf is the fidelity authority).
% Line splits inside nodes follow the accepted edition's render
% exactly. Node ranges carry the single trailing unit
% (\numrange{6}{12}~\unit{\hour}) per the settled FC convention; the
% caption keeps the repeated-unit \qtyrange form.

\begin{tikzpicture}[node distance = 5cm, auto]
  \node [start] (init) {\textarabic{خذ بادئتك}\\ \textarabic{العادية أو}\\ \textarabic{الصلبة}};
  \node [block, right of=init] (feed_new_ratio) {\textarabic{اخلط \LR{1 g}}\\ \textarabic{بادئة موجودة،}\\ \textarabic{و\LR{5 g} دقيق}\\ \textarabic{و\LR{25 g} ماء}};
  \node [decision, below of=feed_new_ratio, node distance=5cm] (ready_signs) {\textarabic{رائحة حامضة}\\ \textarabic{لبنية زبادية}\\ \textarabic{وفقاعات}\\ \textarabic{ظاهرة على}\\ \textarabic{الدقيق؟}};
  \node [block, right of=ready_signs, node distance=4cm] (feed_again) {\textarabic{غذِّ مجددًا}\\ \textarabic{بنسبة}\\ \LR{1:5:25}};
  \node [block, left of=ready_signs, node distance=5cm] (last_feed) {\textarabic{غذِّ مرة}\\ \textarabic{أخيرة}};
  \node [success, below of=last_feed, node distance=4cm] (bread_dough) {\textarabic{اصنع عجين}\\ \textarabic{الخبز}};
  \path [line] (init) -- (feed_new_ratio);
  \path [line] (feed_new_ratio) -- node{\textarabic{انتظر \LR{\qty{24}{\hour}}}} (ready_signs);
  \path [line] (feed_again) -- node[anchor=east]  {} ++(2.2,0) |- (feed_new_ratio);
  \path [line] (ready_signs) -- node{\textarabic{لا}} (feed_again);
  \path [line] (ready_signs) -- node[above=2pt]{\textarabic{نعم}} (last_feed);
  \path [line] (last_feed) -- node{\textarabic{بعد \LR{\numrange{6}{12}~\unit{\hour}}}} (bread_dough);

  \node [above of=feed_again, text width=5em, align=center, node distance=3cm] (repeat_text) {\textarabic{كرر \LR{3} مرات}};
  \draw [line] ($(repeat_text) +(0, 1 cm)$) arc (90:420:1cm);
\end{tikzpicture}

\caption[التحويل إلى بادئة سائلة]{عملية تحويل بادئتك العادية أو الصلبة إلى بادئة سائلة. تستغرق العملية بأكملها نحو \num{3} أيام. وكلما أطلت إبقاء بادئتك عند مستوى الترطيب المقترح، تأقلمت الكائنات الدقيقة أكثر. يُنصح بالاحتفاظ بنسخة احتياطية من بادئتك الأصلية، لأن البيئة السائلة سترشّح الكائنات الدقيقة اللاهوائية. وهذا يعزز البكتيريا المنتجة لحمض اللبن بدلًا من حمض الخليك، فتُستشعر الحمضية الناتجة أخف. وحين تبدأ ببادئة سائلة ستُضفي بادئتك الصلبة نوتات لبنية خفيفة، بينما إن بدأت هذه العملية ببادئة عادية فستحمل البادئة الصلبة الناتجة نوتات لبنية وخلّية معًا.}\label{flc:liquid-starter-conversion}
\end{flowchart}

تُصنع البادئة السائلة بترطيب يقارب \LR{\qty{500}{\percent}}، أي أنها تحوي ماءً يفوق الدقيق بكثير. وتغيّر هذه الطبقة الإضافية من الماء فوق الدقيق الميكروبيوم \textenglish{(microbiome)} في بادئتك.

بإدخال هذه الطبقة من الماء يقل الأكسجين المتاح على مدار التخمير، وبذلك تتوقف بادئتك عن إنتاج حمض الخليك. وتزدهر في هذه البيئة بكتيريا حمض اللبن غير المتجانسة التخمّر \textenglish{(heterofermentative)}. وهذه حيلة رشيقة لتغيير النمط النكهي لبادئتك من الخلّي إلى اللبني؛ فبادئتك ستكتسب نوتات كريمية لبنية. والمثير للاهتمام أنه عند تغيير الترطيب من جديد تحافظ بادئتك على النمط النكهي للبادئة السائلة، لكنها تستفيد ثانيةً من نشاط خميري معزز. وتحويل البادئة إلى سائلة لا رجعة فيه؛ فبتغييرك إلى بادئة سائلة تنتقي بشكل دائم مجموعة فرعية من الكائنات الدقيقة تعمل أفضل في البيئة الأكثر سيولة. لذا حتى بعد العودة إلى بادئة عادية أو صلبة تبقى مجموعة الكائنات الدقيقة التي أحدثها التحويل السائل. ولهذا يُنصح بالاحتفاظ بنسخة احتياطية من البادئة قبل التحويل إلى السائلة.

لبدء التحويل، خذ نحو \LR{\qty{1}{\gram}} من بادئتك وامزجها مع \LR{\qty{5}{\gram}} من الدقيق و\LR{\qty{25}{\gram}} من الماء، ثم اقلب الجميع قلبه جيدًا. بعد بضع دقائق سيبدأ الدقيق بالترسّب نحو قاع الوعاء. كرر هذه العملية على مدى بضعة أيام، ورجّ البادئة برفق لتشاهد إن كنت تلمح فقاعات \LR{\ch{CO2}} دقيقة تتحرك في السائل؛ إنها علامة طيبة على جاهزية البادئة. واستخدم أنفك لشم البادئة؛ ينبغي أن تحمل نوتة كريمية لبنية.

ولأن النشاط البكتيري فيها أعلى، تعمل هذه البادئة أفضل ما يكون مع دقيق قوي جدًا يصمد أمام فترة تخمير طويلة. ولن تنجح مع دقيق قمحي ضعيف. أما إن لم يهمك خبز رغيف قائم بذاته، فيمكنك استخدامها بسهولة مع قالب الخبز. وتشتغل هذه البادئة ببراعة أيضًا عند صنع عجين أقراص مسطّحة غني. ولكي أستخدمها أرجّ وعاء البادئة حتى أرى المكونات وقد تجانس بعضها مع بعض، ثم أستخدم نحو \LR{\qty{5}{\percent}} منها بمنظومة الحساب الخبازي. فلكل \LR{\qty{1000}{\gram}} من الدقيق ذلك نحو \LR{\qty{50}{\gram}} من البادئة السائلة. ولأنها سائلة جدًا عليك إدخال هذه الـ\LR{\qty{50}{\gram}} في حساب سوائلك. وأنا أعامل البادئة مباشرةً كسائل في الوصفات؛ فإن دعت الوصفة إلى \LR{\qty{600}{\gram}} من الماء واستخدمت \LR{\qty{50}{\gram}} من البادئة، أكتفي حينها باستخدام \LR{\qty{550}{\gram}} من الماء فقط.

هذا النوع من البادئة مقاتل ممتاز للعفن أيضًا. فبإبعادك الأكسجين من المعادلة لا يستطيع العفن الهوائي النمو على الوجه الصحيح. وإن كانت بادئتك تعاني مشكلة عفن فقد يكون التحويل إلى السائلة هو العلاج. خذ قطعة من بادئتك تشك في نمو العفن فيها، وطبّق التحويل كما ذكرنا سابقًا. سيُقعّد العفن غالبًا مع نفاد غذائه، وكل تغذية جديدة تقلّل من أبواغ العفن؛ إذ لم تعد الأبواغ تستطيع إعادة التنشيط في الظروف اللاهوائية.

والسائل الطافي فوق بادئتك مورد ممتاز يمكنك استخدامه لصنع الصلصات. فإذا أحببت إضافة بعض الحمضية، صفِّ الجزء السائل من بادئتك واستخدمه. وقد استخدمته مرات لا تُحصى لصنع صلصات حارة مخمَّرة لبنيًا \textenglish{(lacto-fermented)}.

\section{البادئة الصلبة}\label{sec:stiff-starter}

\begin{figure}[!htb]
\centering
\includegraphics[width=0.55\textwidth]{sourdough-starter-stiff}
\caption[البادئة الصلبة مقلوبة]{بادئة خبز حامض صلبة استخدمتها لصنع عجين شتولن \textenglish{(Stollen)} لعيد الميلاد. لاحظ الفقاعات على حافة الوعاء. لا يسقط العجين من الوعاء، لكنه سيهبط في الوعاء حتمًا لحظة انهيار بنية الغلوتين بفعل التخمير.}\label{flc:stiff-sourdough-starter}
\end{figure}

البادئة الصلبة هي الأكثر جفافًا بين جميع البادئات؛ فترطيبها يقارب \LR{\qtyrange{50}{60}{\percent}}. أي أنك لكل \LR{\qty{100}{\gram}} من الدقيق تستخدم نحو \LR{\qtyrange{50}{60}{\gram}} من الماء. وإن عجزت عن خلط الدقيق والماء لجفاف الخليط أكثر من اللازم فعليك زيادة كمية الماء، وهذا كثيرًا ما يحدث عند صنع البادئة بدقيق القمح الكامل أو الجاودار؛ فكلما زاد النخالة في دقيقك زاد ما يستطيع امتصاصه من ماء. ينبغي أن تكون قوام البادئة الصلبة قريبًا من قوام عجين المعكرونة أو البيتزا. وعند خلط البادئة لا ينبغي أن تبقى كتل من الدقيق. جرّب وضع البادئة على سطح مطبخك؛ فحين ترفعها ينبغي أن تلتصق قليلًا بسطح مطبخك، وهذا الاختبار يدل على أنك رطّبت الدقيق ما يكفي. أما إذا كان الخليط جافًا أكثر من اللازم فسيتثاقل معدل التخمير كثيرًا وستبدو البادئة خاملة. ينبغي أن تكون البادئة أجف بكثير من البادئة العادية، لكن دون أن تجف جفافًا مبالغًا فيه. انظر الشكل~\LR{\ref{fig:stiff-starter-dry-check}} لمثال مرئي على مستوى الترطيب المطلوب للبادئة.

\begin{figure}[!htb]
\centering
\includegraphics[width=0.75\textwidth]{stiff-starter-dry-check}
\caption[بادئة صلبة جافة وأخرى مثالية الترطيب]{صورة تُريك بادئة صلبة جافة أكثر من اللازم وأخرى مثالية الترطيب. لا ينبغي أن تحوي البادئة كتلًا من الدقيق، وأن تلتصق قليلًا بسطح مطبخك. البادئة في الصورة مصنوعة من دقيق القمح الكامل.}\label{fig:stiff-starter-dry-check}
\end{figure}

\begin{flowchart}[!htb]
\centering
% fig-stiff-starter-conversion-ar.tex — Arabic edition of the stiff
% starter conversion flowchart (ch. 6, المخطط الانسيابي 6.2)
%
% Original: The Sourdough Framework by Hendrik Kleinwächter,
% the-sourdough-framework.com, CC BY-SA 4.0. The Arabic edition inherits
% CC BY-SA 4.0.
%
% R1-FIGURES route: re-derived from the upstream post-#492 source
% figures/fig-stiff-starter-conversion.tex (node styles from
% figures/flowcharts_tikz.tex and the node/edge geometry verbatim,
% including the at (feed_new_ratio |- next_day) and
% at (ready_signs -| add_water) placements, the
% ++(0, 1.3) -| (too_dry.north) loop, the (feed_again |- last_feed)
% -| (feed_again.south) return edge and the (90:-210:1cm) repeat arc),
% with the label text recovered from the accepted edition's figure
% renders (label recovery, not retranslation; the accepted edition
% book-ar-fixed.pdf is the fidelity authority). Line splits inside
% nodes follow the accepted edition's render exactly. Node ranges
% carry the single trailing unit (\numrange{6}{12}~\unit{\hour}) per
% the settled FC convention; the caption keeps the repeated-unit
% \qtyrange form.

\begin{tikzpicture}[node distance = 4cm, auto]
  \node [start] (init) {\textarabic{خذ بادئتك}\\ \textarabic{العادية أو}\\ \textarabic{السائلة}};
  \node [block, right of=init, node distance = 4cm] (feed_new_ratio) {\textarabic{اخلط \LR{10 g}}\\ \textarabic{بادئة موجودة،}\\ \textarabic{و\LR{50 g} دقيق}\\ \textarabic{و\LR{25 g} ماء}};
  \node [decision, right of=feed_new_ratio, node distance=5cm] (too_dry) {\textarabic{البادئة جافة}\\ \textarabic{جدًا ويصعب}\\ \textarabic{خلطها؟}};
  \node [block, right of=too_dry, node distance=4cm] (add_water) {\textarabic{أضف المزيد}\\ \textarabic{من الماء}};
  \node [block, below of=too_dry] (next_day) {\textarabic{انتظر}\\ \LR{\qty{24}{\hour}}};
  \node [block] at (feed_new_ratio |- next_day) (feed_again) {\textarabic{غذِّ مجددًا}\\ \textarabic{بنسبة}\\ \LR{1:5:2.5}};
  \node [decision, below of=next_day, node distance=3.5cm] (ready_signs) {\textarabic{زيادة في الحجم}\\ \textarabic{ورائحة}\\ \textarabic{حامضة؟}};
  \node [block] at (ready_signs -| add_water) (last_feed) {\textarabic{غذِّ مرة}\\ \textarabic{أخيرة}};
  \node [success, below of=last_feed, node distance=3cm] (bread_dough) {\textarabic{اصنع عجين}\\ \textarabic{الخبز}};
  \path [line] (init) -- (feed_new_ratio);
  \path [line] (feed_again) -- (feed_new_ratio);
  \path [line] (next_day) -- (ready_signs);
  \path [line] (ready_signs) -- node{\textarabic{لا}} (feed_again |- last_feed) -| (feed_again.south);
  \path [line] (ready_signs) -- node{\textarabic{نعم}} (last_feed);
  \path [line] (last_feed) -- node{\textarabic{بعد \LR{\numrange{6}{12}~\unit{\hour}}}} (bread_dough);
  \path [line] (feed_new_ratio) -- (too_dry);
  \path [line] (add_water.north) -- node{} ++(0, 1.3) -| (too_dry.north);
  \path [line] (too_dry) -- node{\textarabic{لا}} (next_day);
  \path [line] (too_dry) -- node{\textarabic{نعم}} (add_water);

  \node [text width=5em, align=center] (repeat_text) at ($(feed_new_ratio)!0.5!(next_day)$) {\textarabic{كرر \LR{3} مرات}};
  \draw [line] ($(repeat_text) +(0, 1 cm)$) arc (90:-210:1cm);
\end{tikzpicture}

\caption[التحويل إلى بادئة صلبة]{عملية تحويل بادئتك العادية إلى بادئة صلبة. تستغرق العملية بأكملها نحو \num{3} أيام. وكلما أطلت إبقاء بادئتك عند مستوى الترطيب المقترح، تأقلمت الكائنات الدقيقة أكثر. تعزز البادئة الصلبة نشاط الخميرة في بادئة الخبز الحامض. يستخدم هذا الدليل مستوى ترطيب \LR{\qty{50}{\percent}} للبادئة، وإن كان العجين صلبًا أكثر من اللازم ففكر في رفعه إلى \LR{\qty{60}{\percent}}.}\label{fig:stiff-starter-conversion}
\end{flowchart}

في البيئة الأصلب تزدهر الخميرة أكثر، أي أنك ستحصل على إنتاج أكبر من \LR{\ch{CO2}} وإنتاج أقل من الحمضية. وفي تجاربي كان هذا نقلة نوعية، لا سيما إن كنت تستخدم دقيقًا أضعف غلوتينًا. فدقيق القمح في بلدي ألمانيا يميل إلى انخفاض الغلوتين. ويُفضَّل لبناء الغلوتين في القمح ظروف دافئة~\cite{gluten+development+temperatures}. وعند اتباع وصفات خبازين آخرين لم أستطع يومًا بلوغ نتائج مماثلة؛ فحين أتبع الأزمان كان عجيني ينهار ببساطة ويصبح لزجًا للغاية. ولم تتغير نتائجي إلا حين بدأت بشراء دقيق قمح أغلى ثمنًا. ولأن ليس الجميع قادرًا على تحمل كلفة هذا الدقيق الخاص ولضآلة توفّره، صادفتُ البادئة الصلبة. أجريت عدة تجارب استخدمت فيها الكمية نفسها من البادئة والدقيق، ولم أغير سوى الترطيب بين جميع البادئات، ثم وضعت بالونًا فوق كل وعاء من الأوعية. وقد انتفخ وعاء البادئة الصلبة انتفاخًا أوضح بكثير، وتبعته البادئة العادية في المركز الثاني، وجاءت البادئة السائلة في المركز الثالث بإنتاج \LR{\ch{CO2}} أقل بكثير.

\begin{figure}[!htb]
\centering
\includegraphics[width=\textwidth]{stollen}
\caption[شتولن عيد الميلاد]{شتولن ألماني لعيد الميلاد مصنوع ببادئة صلبة بدلًا من الخميرة.}\label{fig:stollen}
\end{figure}

ثم اشتريت دقيق كيك رخيصًا منخفض الغلوتين من السوبرماركت المجاور، وهو الدقيق الذي سبق لي صداعًا هائلًا في الماضي. صنعت خبزًا حامضًا تمامًا بالطريقة التي أعتادها، وقد اضطررت إلى تخفيض الترطيب قليلًا لأن الدقيق منخفض الغلوتين لا يمتص الماء كثيرًا. ثم استبدلت البادئة بالبادئة الصلبة. فأحسّ العجين برائع، وصار فجأة قادرًا على الصمود أمام فترة تخمير أطول بكثير. وكان للخبز انتفاخ ممتاز في الفرن ومذاق لطيف للغاية. وما زلت أبحث عن تفسير علمي وافٍ لسبب نشاط جزء الخميرة في العجين أكثر. ولعلها ليست أكثر فعلًا؛ فقد يكون البكتيريا هو المثبَّت بفعل شح الماء.

وعند صنع البادئة الصلبة ابدأ باستخدام نحو \LR{\qty{50}{\percent}} من الماء. وإن كنت تستخدم دقيق قمح كامل، أو دقيقًا قويًا، ففكر في الارتقاء إلى \LR{\qty{60}{\percent}}. ينبغي أن تمتزج جميع المكونات امتزاجًا ممتازًا ولا تبقى أي كتل دقيق. وهذا خطأ شائع رأيته عند محاولة الناس صنع البادئة الصلبة: نعم ينبغي أن تكون جافة، لكن ليس إلى حد كتلة إسمنتية. وإن سبق أن صنعت عجين معكرونة، فيجب أن يحسّ هذا العجين بالضبط بالمثل.

وللتحقق من جاهزية بادئتك الصلبة ابحث عن قبة منتفخة. ابحث أيضًا عن جيوب هواء على جوانب وعائك. واستخدم أنفك لشم البادئة؛ ينبغي أن تكون رائحتها لطيفة، كما أنها تميل إلى أن تكون أكثر كحولية بكثير من البادئات الأخرى.

وعند استخدام بادئة صلبة، اعتمد نحو \LR{\qtyrange{1}{20}{\percent}} منها بمنظومة الحساب الخبازي بالنسبة إلى عجينك، وذلك بحسب نضج بادئتك. ففي الصيف أستخدم عادة نحو \LR{\qtyrange{1}{10}{\percent}} وفي الشتاء نحو \LR{\qty{20}{\percent}}. وبهذه الطريقة تستطيع أيضًا ضبط سرعة التخمير. فإن كان الجو حارًا جدًا حيث تعيش ففكر في خفض كمية البادئة إلى \LR{\qtyrange{1}{5}{\percent}}، وإن كان باردًا جدًا في منطقتك ففكر في رفع كمية البادئة حتى \LR{\qty{30}{\percent}}. وقد يكون خلط البادئة الصلبة مزعجًا قليلًا لأنها بالكاد تتجانس مع بقية العجين. في هذه الحالة يمكنك محاولة إذابة البادئة في الماء الذي ستصنع به عجينك؛ سيسهل ذلك الخلط كثيرًا.

\section{الليفيتو مادر أو الباستا مادر}

يصنّف الليفيتو مادر \textenglish{(lievito madre)}، المعروف أيضًا بالباستا مادر \textenglish{(pasta madre)}، في الصف نفسه الذي تصنّف فيه البادئة الصلبة. وبعد ساعات من البحث لم أجد فرقًا بين الباستا مادر والليفيتو مادر؛ فالمصطلحان يبدو أنهما يستخدمان بالتبادل في المراجع.

في كثير من الوصفات تُصنع هذه البادئة مباشرةً من الفواكه المجففة أو الطازجة، ويمكنك أيضًا صنع بادئة من أوراق من حديقتك. وكما وصفنا سابقًا، تستهلك خميرة وبكتيريا البرية الغلوكوز من أوراق النبات. وكل الخيارات تنفع. وعند صنع بادئة من فواكه مجففة مباشرةً قد يخلو الأمر أحيانًا من الجزء البكتيري من التخمير. والحمضية مهمة جدًا لتطهير بادئتك من مسببات الأمراض المحتملة. فإذا قررت صنع بادئتك من فواكه فتأكد من أنها تحمض أيضًا على الوجه الصحيح عند صنع العجين. وأداة مثل مقياس الرقم الهيدروجيني \textenglish{(pH meter)} قد تساعدك على أفضل وجه. وبشكل عام، كلما انخفض الرقم الهيدروجيني ارتفعت الحمضية. ويجب أن تبقى الحمضية تحت \num{4.2} لتعرف أن بادئتك تنتج حمضية كافية.

ينظف بعض الخبازين الليفيتو مادر في حمّام من الماء، على افتراض أن ذلك يزيل الحمضية الزائدة. وفي تجاربي الخاصة لم أستطع تأكيد هذه المنهجية؛ فالحمضية بقيت على حالها. والسبب الوحيد الذي قد يجعل هذا الإجراء ذا معنى هو أن تحاول في الوقت نفسه تعزيز الكائنات الدقيقة اللاهوائية. غير أن البادئة حينها ستحتاج إلى البقاء في هذه البيئة مدة طويلة حقًا، لا بضع ساعات فقط.

\section{الخاتمة}\label{sec:starter-type-conclusion}

الخبز بالخبز الحامض أمر بسيط؛ إنه مجرد دقيق وماء. وحين ترى وصفة خباز خبير تتساءل: هل هذا هو كل شيء؟ لا شيء سواه؟ وأشعر أن هذا قد يكون السبب وراء تعقيد بعض الخبازين لإجراءات تغذية بادئاتهم، فيلجؤون إلى عدة تغذيات يوميًا بنسبة محددة. هذا يمنح الخباز إحساسًا بالنخبوية قليلًا. وطبعًا مع الوقت، مع اتباع المزيد من الناس هذا الإجراء، أصبح نبوءة تحقق نفسها: فكلما ازددت خبرة ازدادت فرص أن يكافئك دليل تغذية بادئة مفتعل بنتائج جميلة. لكن السبب ليس في روتين البادئة، بل في أنك تفهم التخمير أفضل وتصبح أقدر على قراءة علامات عجينك.

وإن اضطررت إلى اختيار نوع واحد من البادئات لاخترت البادئة الصلبة؛ ففي حالات كثيرة ستحقق لك نتائج رائعة متسقةً مع جهد ضئيل. وفي تجربتي يمكنك صنع أي عجين خميري واستبدال الخميرة مباشرةً بالبادئة الصلبة، بل ستحقق نتائج أفضل حتى مع البادئة الصلبة.

وأخيرًا، أيًا كان نوع البادئة الذي تختاره، يمكنك التحكم في درجة حمضية عجينك كما تريد. فكلما أطلت التخمير تراكمت الحمضية أكثر. والفرق الوحيد أنه عند زيادة حجم معينة تنتج البادئة الصلبة أقل الحمضية. فعند زيادة حجم قدرها \LR{\qty{100}{\percent}} تكون البادئة السائلة قد أنتجت أكثر الحمضية، تليها البادئة العادية ثم البادئة الصلبة. وإن انتظرت وقتًا كافيًا ستكون البادئة الصلبة قد أنتجت القدر نفسه من الحمضية الذي أنتجته البادئات الأخرى، لكنها قبل ذلك تكون قد أنتجت كثيرًا من \LR{\ch{CO2}} أيضًا. فإذا أحببت النكهة الحامضة فعليك إطالة التخمير، وهذا يعني أيضًا أنك ستحتاج إما إلى الخبز في قالب الخبز أو إلى دقيق غلوتيني قوي جدًا يصمد أمام أزمنة تخمير طويلة.
